The ice sang before it broke.

Tormund Giantsbane had heard the Wall make many sounds in his life — the low groan of it settling on cold nights, the drip and crack of thaw in a warm summer, the wind keening through the tunnels the free folk had cut generations before he was born to slip south and back again. He had heard it weep. He had never heard it sing.

It began as a hum, so faint he took it at first for the blood in his own ears. He was standing atop the ice at Eastwatch-by-the-Sea, three hundred feet above the black water where it met the black rock, and the hum came up through his boots before it ever reached his ears at all. Dolorous Edd felt it too. Tormund saw him go still beside the brazier, head cocked like a dog that has heard its master’s step on a stair the master does not know is there.

“You feel that,” Edd said. It was not a question. Edd rarely asked questions. He seemed to have decided some years back that the world’s answers were uniformly disappointing and had stopped soliciting them.

“Aye.”

“I don’t like it.”

“You don’t like anything.”

“I like some things,” said Edd. “I liked it better before I felt that.”

Below them, beyond the black water, beyond the shapes of the drowned and reborn dead that had followed the Night King’s host up out of the haunted forest and along the frozen shore, something vast was moving. Tormund had seen the dead walk. He had seen them run, when the fires were lit and the wights came on screaming without voices, mouths open on nothing. He had never seen the sea itself seem to flinch, and that was what he saw now — the black water shivering all along the shoreline, the way a horse’s hide shivers when a fly lands on it, as though the cold itself were a living thing drawing back from something colder still.

Then the light came, and the singing stopped.

It rose out of the dark somewhere north and east of the ice, low over the tree line at first, and Tormund’s first thought — the thought of a man’s mind reaching for the nearest explanation his life has given him, before it reaches for the true one — was fire. Some vast, unnatural torch, or a comet fallen wrong, burning blue instead of the orange a comet ought to burn. It moved too smoothly for fire, though, too unhurried, tracking a slow arc against the stars that had nothing of a flame’s flicker in it, and it made no sound at all, which unsettled him worse than any roar could have, because everything Tormund Giantsbane had ever feared in his life had announced itself with noise, and this came on silent as a held breath.

“Is that a dragon?” Edd’s voice had gone strange, higher than Tormund had ever heard it. “Tell me that’s not a dragon, Tormund. Tell me the Dragon Queen’s flown one north to help us and forgot to send word first.”

Tormund’s stomach had already answered that question before his mouth caught up to it. He had stood on ice beyond this same Wall not half a year past and watched the Night King’s spear take one of the silver queen’s three children out of the sky, watched it sink beneath black water and not come up again — or not come up again \textit{as itself.} He knew, some old cold certainty settling into him even as some other, more foolish part of him reached for hope anyway, that whatever was banking low over the tree line now was not a living queen’s dragon flown north out of mercy. It was the one he’d already watched die once, wearing its death the way a man wears a wound he never let heal, and some small, treacherous corner of him spent one last heartbeat hoping he was wrong about which one it was, before his eyes finished giving him the truth of it regardless.

It was not blue like a jewel, the way the singers claimed the eyes of the dead things were blue. It was blue the way a wound is blue before it goes black — a color with sickness in it, glowing faint along the ridges of wings that did not move the way a living wing moved, stiff in the joints, laboring, as though whatever drove it forward cared nothing for how much the effort cost the body doing the flying. And when it opened its jaws over the ice, no roar came out of it, no living dragon’s fury — only that same silent, patient blue light, pouring down onto the Wall the way a man pours water down a hole he already knows the shape of, unhurried, because why would you hurry, when you had waited eight thousand years and the only thing left between you and the world was seven hundred feet of ice that had never once in all that span been asked to hold against dragonfire, because no one, in eight thousand years, had thought that was a thing the Wall would ever have to do.

“That’s no living queen’s dragon,” Tormund said, and understood, saying it, that he had known from the first silent heartbeat and had simply needed the saying of it to stop pretending otherwise. “That’s the one we already lost. Wearing its death like a coat that still fits.”

“Get the horn,” Edd said, and there was no dolor in his voice now, none at all. “Get the bloody horn, Tormund, get it now—“

Tormund did not remember afterward whether he had sounded it, or whether someone else had, or whether the sound he heard was the horn at all and not simply the Wall itself giving voice to what was being done to it. What he remembered was the light coming down not on the men atop the ice, not first, but on the ice beneath their feet — a long lance of that sick blue fire, drawn down the face of the Wall like a man drawing a blade down a whetstone, patient, testing, as though whatever rode that dying beast wanted to be certain, before it spent the effort, that this was a wound the Wall would not simply heal itself of by morning.

It held for the space of three breaths. Tormund counted them without meaning to, the way a man counts his heartbeats when a wound has gone too deep to feel yet. One. The ice groaning, a sound so low he felt it more than heard it. Two. A crack opening somewhere far below, a black seam spreading up through white ice like a river finding a course, and men — his men, Edd’s men, black brothers he had shared fires and jokes and hard words with for two years now, men whose names he would still be saying to himself years later when he could not sleep — men beginning to run, though there was nowhere on the Wall itself to run to. Three. The world tilted.

He did not remember deciding to jump. He remembered the ice going out from under him the way a wave goes out from under a swimmer, not falling so much as being \textit{removed}, and he remembered grabbing for Edd’s collar and missing, and he remembered — this above all, this was the thing that would come to him unbidden at strange hours for the rest of his life — Edd’s face in that last instant, not afraid, or not only afraid. Almost amused. As if some old jape had finally landed, years too late to laugh at properly.

Then he was in the water, and the water was the temperature of dying, and the sky above him was full of falling ice the size of holdfasts, and somewhere behind him and above him and in every direction at once the Wall of Westeros, seven hundred feet of ice and eight thousand years of purpose, came down into the sea in a white roar that swallowed every other sound in the world, including, mercifully, for a little while, his own.

\sceneBreak

He did not die. He would turn that fact over in his hands for the rest of his days like a stone he could not put down, uncertain whether it was a mercy or an oversight.

A fishing skiff out of Eastwatch’s harbor — one of the few boats that had not been crushed by the falling ice or swamped by the wave that followed — found him half a mile out, clinging to a spar of black wood that might once have been a support beam or might once have been a man. They pulled him aboard more dead than not, and it was two more days before he could speak, and when he could speak the first thing he said, to a fisherman’s wife who had been spooning broth into him like he was an infant, was: “The Watch.”

“There’s no Watch,” she told him, not unkindly. “There’s no Wall, ser. Just the Wall’s own self, is all, floating out there. Ice islands big as this boat, some of them. Bigger.”

“Edd,” he said.

She had not known an Edd. She had known no one at Eastwatch save what the sea had given her, and the sea, in the end, had given her only Tormund Giantsbane, one man out of some hundreds, and even he was not certain, lying in that fisherman’s hut with his lungs still full of salt, whether he ought to be grateful.

It was a passing trader, three days later, who first said the name \textit{Viserion} to him plainly, having heard it already from some other survivor further down the coast, and Tormund found, hearing it said aloud by another man’s mouth, that it settled into him strangely — not as news, since he’d known well enough what he’d watched die twice now, but as confirmation that the rest of the world was going to have to learn to say that name the way he already had, with the particular dread reserved for a wonder that had been made into a weapon. He had hoped, some small foolish corner of him, that he might have been wrong in the dark and the panic. He was not wrong. Whatever had broken the Wall had been hers once, before the Night King unmade what a dragon was and left only the shape of one walking around in it still, and Tormund did not know which was worse to carry south with him — that dragons could die at all, or that even in dying they weren’t done being made to matter.

He set out south three days after that on a stolen horse he did not feel very guilty about stealing, because the man he stole it from was already three days dead of something that was not the cold and not the dead things either, just an old man’s heart giving out at the sight of a floating Wall, and Tormund reasoned the horse had no more use for its owner’s blessing than the owner did. He rode hard, and he rode alone, because there was no one left at Eastwatch to ride with him, and all the way south the sky behind him stayed the color of a healing bruise, day and night both, as though something enormous had been struck and had not yet decided whether it meant to fall down.

He did not know, riding, whether the Night King had come through the breach yet, or whether he waited, the way a man might wait at a broken gate to be sure no arrow answered him. He did not know how many had followed the ice down, how many wights had gone under and how many had simply walked across, the sea itself no longer proof against them now that the thing that had held it back for eight thousand years was so much broken glass floating in black water.

He knew only that someone needed telling, and that the someone needed telling \textit{soon}, and that he was, as far as he could presently determine, the only man breathing who could tell them.

He turned his stolen horse’s head toward Winterfell, and rode.